%=====================================================================
% CARGAS DEDICADAS - REPRESENTACAO REVISADA
%=====================================================================
\section{Cargas dedicadas e topologias de alimentação}

\begin{frame}{Mesmo valor de tensão, topologias diferentes}
\small
\begin{columns}[T,onlytextwidth]
\column{.49\textwidth}
\begin{caixaProjeto}[title=Fase--neutro]
\centering
\begin{tikzpicture}[scale=.72,every node/.style={font=\scriptsize,align=center}]
\node[draw=corPrim,fill=corDef!5,rounded corners,minimum width=2cm,minimum height=.8cm](q)at(-2,0){origem};
\node[draw=corAccent!85!black,fill=corAccent!7,rounded corners,minimum width=2cm,minimum height=.8cm](c)at(2,0){carga};
\draw[thick,corPrim]($(q.east)+(0,.18)$)--($(c.west)+(0,.18)$)node[midway,above]{F};
\draw[thick,corDef]($(q.east)+(0,-.18)$)--($(c.west)+(0,-.18)$)node[midway,below]{N};
\end{tikzpicture}
\end{caixaProjeto}
\column{.49\textwidth}
\begin{caixaProjeto}[title=Fase--fase]
\centering
\begin{tikzpicture}[scale=.72,every node/.style={font=\scriptsize,align=center}]
\node[draw=corPrim,fill=corDef!5,rounded corners,minimum width=2cm,minimum height=.8cm](q)at(-2,0){origem};
\node[draw=corAccent!85!black,fill=corAccent!7,rounded corners,minimum width=2cm,minimum height=.8cm](c)at(2,0){carga};
\draw[thick,corPrim]($(q.east)+(0,.18)$)--($(c.west)+(0,.18)$)node[midway,above]{F1};
\draw[thick,corNorma]($(q.east)+(0,-.18)$)--($(c.west)+(0,-.18)$)node[midway,below]{F2};
\end{tikzpicture}
\end{caixaProjeto}
\end{columns}
\begin{caixaObs}
O valor nominal de tensão não informa sozinho quais condutores compõem o circuito. Essa diferença precisa aparecer no diagrama.
\end{caixaObs}
\end{frame}

\begin{frame}{Carga dedicada — mostrar a cadeia física completa}
\centering
\begin{tikzpicture}[scale=.82,every node/.style={font=\scriptsize,align=center}]
\node[draw=corPrim,fill=corDef!5,rounded corners,minimum width=2.1cm,minimum height=.85cm](q)at(-5.0,0){quadro};
\node[draw=corNorma,fill=corNorma!5,rounded corners,minimum width=2.1cm,minimum height=.85cm](p)at(-1.7,0){proteção /\\seccionamento};
\node[draw=corDef,fill=corDef!6,rounded corners,minimum width=2.1cm,minimum height=.85cm](x)at(1.7,0){conexão};
\node[draw=corAccent!85!black,fill=corAccent!7,rounded corners,minimum width=2.1cm,minimum height=.85cm](c)at(5.0,0){carga};
\draw[thick,corPrim,-{Stealth}](q)--(p)--(x)--(c);
\draw[thick,corNorma](-5,-1.35)--(5,-1.35)--(c.south);
\node[font=\tiny]at(0,-1.65){condutor de proteção representado separadamente};
\end{tikzpicture}
\begin{caixaProjeto}[title=Melhoria gráfica]
A representação mostra origem, proteção, ponto de conexão e equipamento como elementos distintos. Isso evita que um único bloco esconda decisões de projeto e interfaces físicas.
\end{caixaProjeto}
\end{frame}

\begin{frame}{Bomba comandada — potência e comando em diagramas coordenados}
\small
\begin{columns}[T,onlytextwidth]
\column{.49\textwidth}
\centering\textbf{Potência}\par\vspace{1mm}
\begin{tikzpicture}[scale=.72,every node/.style={font=\scriptsize,align=center}]
\node[draw=corPrim,fill=corDef!5,minimum width=2cm](q)at(0,2.4){proteção};
\node[draw=corNorma,fill=corNorma!5,minimum width=2cm](k)at(0,1.1){K1};
\node[draw=corAccent!85!black,fill=corAccent!7,circle,minimum size=.95cm](m)at(0,-.5){M};
\draw[thick,corPrim,-{Stealth}](q)--(k)--(m);
\end{tikzpicture}
\column{.49\textwidth}
\centering\textbf{Comando}\par\vspace{1mm}
\begin{tikzpicture}[scale=.72,every node/.style={font=\scriptsize,align=center}]
\node[draw=corPrim,fill=corDef!5,minimum width=2cm](s)at(-2.2,1.3){sensor / boia};
\node[draw=corNorma,fill=corNorma!5,minimum width=2cm](b)at(2.2,1.3){bobina K1};
\draw[thick,corPrim,-{Stealth}](s)--(b);
\draw[dashed,corAccent,-{Stealth}](b.south)--(2.2,-.4)node[below]{atua em K1};
\end{tikzpicture}
\end{columns}
\begin{caixaObs}
A mesma referência K1 liga os dois desenhos. Essa referência cruzada é mais útil ao aluno do que uma única seta entre sensor e motor.
\end{caixaObs}
\end{frame}
